\documentclass[11pt,oneside]{book}
\usepackage[a4paper,margin=28mm]{geometry}
\usepackage{fontspec}
\defaultfontfeatures{Ligatures=TeX}
\setmainfont{DejaVu Serif}
\setsansfont{DejaVu Sans}
\setmonofont{DejaVu Sans Mono}
\usepackage{unicode-math}
\setmathfont{DejaVu Math TeX Gyre}
\usepackage{microtype}
\setlength{\emergencystretch}{3em}
\usepackage{amsmath,amsthm,mathtools}
\usepackage{booktabs,longtable,array}
\usepackage{graphicx}
\graphicspath{{../figures/}}
\usepackage{enumitem}
\usepackage{csquotes}
\usepackage{hyperref}
\usepackage{xcolor}
\usepackage{fancyhdr}
\usepackage{listings}
\usepackage{caption}
\usepackage{subcaption}
\usepackage{setspace}
\usepackage{imakeidx}
\makeindex

\hypersetup{
  pdftitle={The Information-Spinor Critical Axis},
  pdfauthor={Adrian Lipa},
  pdfsubject={Shannon balance, Euler-Berry holonomy, and Riemann zeta symmetry},
  pdfkeywords={Riemann hypothesis, zeta function, Shannon entropy, Berry phase, spinor, Bloch sphere, Hilbert-Polya},
  colorlinks=true,
  linkcolor=blue!45!black,
  citecolor=blue!45!black,
  urlcolor=blue!45!black
}

\pagestyle{fancy}
\fancyhf{}
\fancyhead[L]{The Information-Spinor Critical Axis}
\fancyhead[R]{Adrian Lipa}
\fancyfoot[C]{\thepage}
\setlength{\headheight}{14pt}

\newtheorem{theorem}{Theorem}[chapter]
\newtheorem{proposition}[theorem]{Proposition}
\newtheorem{lemma}[theorem]{Lemma}
\newtheorem{corollary}[theorem]{Corollary}
\newtheorem{definition}[theorem]{Definition}
\newtheorem{conjecture}[theorem]{Conjecture}
\newtheorem{assumption}[theorem]{Assumption}
\newtheorem{openproblem}[theorem]{Open Problem}
\theoremstyle{remark}
\newtheorem{remark}[theorem]{Remark}
\newtheorem{warning}[theorem]{Epistemic Warning}

\newcommand{\Ree}{\operatorname{Re}}
\newcommand{\Imm}{\operatorname{Im}}
\newcommand{\dd}{\,\mathrm{d}}
\newcommand{\ii}{\mathrm{i}}
\newcommand{\ee}{\mathrm{e}}
\newcommand{\Hbin}{H_{\mathrm{bin}}}
\newcommand{\ket}[1]{\lvert #1\rangle}
\newcommand{\bra}[1]{\langle #1\rvert}
\newcommand{\braket}[2]{\langle #1\mid #2\rangle}
\newcommand{\status}[1]{\textsf{#1}}

\begin{document}

\frontmatter
\begin{titlepage}
\centering
\vspace*{2.0cm}
{\Huge\bfseries The Information-Spinor Critical Axis\par}
\vspace{0.5cm}
{\Large Shannon Balance, Euler-Berry Holonomy, and the Riemann Zeta Symmetry\par}
\vspace{1.5cm}
{\Large Adrian Lipa\par}
\vspace{0.3cm}
{\large Independent Researcher, Doncaster, United Kingdom\par}
\vfill
{\large Research Monograph v0.1\par}
{\large 29 July 2026\par}
\vspace{1cm}
\begin{minipage}{0.88\textwidth}
\small
\textbf{Status statement.} This monograph proves a set of exact mathematical bridges among binary Shannon balance, complementary two-channel interference, qubit/Bloch geometry, Berry holonomy, and the fixed axis of the completed Riemann zeta functional equation. It then formulates a conditional critical-axis selection theorem. It does \emph{not} claim an unconditional proof of the Riemann hypothesis. The remaining implication is isolated as an explicit operator/kernel problem.
\end{minipage}
\end{titlepage}

\chapter*{Abstract}
\addcontentsline{toc}{chapter}{Abstract}
The critical line $\Ree s=1/2$ simultaneously appears as the fixed set of the zeta involution, the maximum-entropy point of a binary probability simplex, the equator of a qubit Bloch sphere, and the unique point at which two complementary amplitudes can cancel under a relative phase of $\pi$. This monograph develops those coincidences into a precise framework. First, it derives the Shannon balance theorem $\Hbin(1/2)=\ln 2$. Second, it embeds the Bernoulli parameter $\sigma$ into the normalized state $\ket{\psi}=\sqrt{\sigma}\ket{0}+e^{i\phi}\sqrt{1-\sigma}\ket{1}$ and proves that exact destructive interference occurs if and only if $\sigma=1/2$ and $\phi=\pi$ modulo $2\pi$. Third, it computes the Berry phase of a full azimuthal cycle and shows that at $\sigma=1/2$ the holonomy is $-1$, the spinorial sign associated with a $2\pi$ cycle. Fourth, it places this geometry beside the completed zeta symmetry $\xi(s)=\xi(1-s)$ and the anti-linear involution $J(s)=1-\overline{s}$, whose fixed set is precisely $\Ree s=1/2$. The Dirichlet eta identity $\eta(s)=(1-2^{1-s})\zeta(s)$ and $\eta(1)=\ln2$ supply an exact parity/alternation bridge, while the Hardy $Z$ function supplies the established critical-line phase normalization. The principal result is conditional: if each nontrivial zero is represented by a locally normalized complementary two-branch cancellation state compatible with the zeta involution, then its real part must equal $1/2$. The unresolved task is therefore not the selection of the axis but the construction of a canonical representation or self-adjoint operator that forces every zero into that local cancellation class.

\chapter*{Claim Ledger}
\addcontentsline{toc}{chapter}{Claim Ledger}
\begin{longtable}{@{}p{0.28\textwidth}p{0.18\textwidth}p{0.46\textwidth}@{}}
\toprule
Claim & Status & Meaning\\
\midrule
Binary entropy has a unique maximum at $\sigma=1/2$ with value $\ln2$. & \status{THEOREM} & Standard calculus result, proved here.\\
The complement involution $\sigma\mapsto1-\sigma$ has the unique fixed point $1/2$. & \status{THEOREM} & Elementary and exact.\\
Two complementary amplitudes cancel exactly at relative phase $\pi$ iff $\sigma=1/2$. & \status{THEOREM} & Exact algebraic result, proved here.\\
The equatorial qubit cycle has Berry holonomy $-1$. & \status{THEOREM} & Exact within the stated gauge convention; the holonomy is gauge invariant.\\
The completed zeta function is symmetric under $s\mapsto1-s$, and its anti-linear fixed axis is $\Ree s=1/2$. & \status{THEOREM} & Classical analytic number theory.\\
$\eta(s)=(1-2^{1-s})\zeta(s)$ and $\eta(1)=\ln2$. & \status{THEOREM} & Classical identity.\\
The parameter identification $\sigma=\Ree s$ is canonical. & \status{MODEL POSTULATE} & Motivated by complement symmetry, not supplied by standard zeta theory.\\
Every nontrivial zeta zero is a local balanced two-channel cancellation state. & \status{OPEN ASSUMPTION} & This is the missing bridge.\\
All nontrivial zeta zeros lie on $\Ree s=1/2$. & \status{OPEN} & The Riemann hypothesis is not proved here.\\
\bottomrule
\end{longtable}

\tableofcontents

\mainmatter
\chapter{Problem, Scope, and Method}
\label{ch:problem}

\section{The proposed mechanism}
The working hypothesis of this monograph can be stated in one sentence:

\begin{quote}
The real coordinate $\sigma=\Ree s$ of a nontrivial zeta zero is selected by the simultaneous requirements of complementary information balance and a phase reversal, so that the unique admissible cancellation axis is $\sigma=1/2$.
\end{quote}

The statement contains three distinct mathematical objects that must not be conflated:
\begin{enumerate}[label=(\roman*)]
\item the analytic symmetry of the completed Riemann zeta function;
\item the binary information geometry of a pair of complementary weights $\sigma$ and $1-\sigma$;
\item the geometry and holonomy of a two-level complex state.
\end{enumerate}
The goal is to derive each layer independently and then identify the precise extra assumption required to connect them.

The central discipline of the project is that a compelling structural coincidence is not automatically a proof. The line $\Ree s=1/2$ is already distinguished by the functional equation. The difficult part of the Riemann hypothesis is to exclude zeros that occur in symmetric pairs away from that line. Any successful argument must therefore do more than rediscover the symmetry. It must supply a positivity, self-adjointness, local fixed-point, or equivalent mechanism that forbids off-axis quartets.

\section{Four levels of assertion}
We use the following hierarchy throughout.

\begin{description}[style=nextline]
\item[THEOREM] A statement proved from explicit assumptions by standard mathematics.
\item[PROPOSITION] A derived statement of narrower scope, often used as a bridge between chapters.
\item[MODEL POSTULATE] An identification introduced to build the information-spinor representation. It is mathematically consistent but not forced by standard zeta theory.
\item[OPEN GAP] A missing implication whose proof would be needed before any claim of a proof of the Riemann hypothesis.
\end{description}

This hierarchy is not cosmetic. It prevents the most common failure mode in speculative work on the zeta function: replacing the sentence ``the critical line is the symmetry axis'' with the much stronger and unproved sentence ``all zeros must lie on the symmetry axis.'' The first is classical. The second is the Riemann hypothesis.

\section{The three appearances of one half}
The number $1/2$ enters the construction in three exact ways.

First, the binary Shannon entropy in natural units,
\begin{equation}
\Hbin(\sigma)=-\sigma\ln\sigma-(1-\sigma)\ln(1-\sigma),
\end{equation}
has a unique maximum at $\sigma=1/2$, with value $\ln2$. This is the point at which the two alternatives carry equal probability weight.

Second, a normalized two-level state may be written as
\begin{equation}
\ket{\psi(\sigma,\phi)}
=\sqrt{\sigma}\ket{0}+e^{i\phi}\sqrt{1-\sigma}\ket{1}.
\end{equation}
At the relative phase $\phi=\pi$, the scalar interference amplitude
\begin{equation}
\mathcal A(\sigma,\pi)=\sqrt{\sigma}-\sqrt{1-\sigma}
\end{equation}
vanishes if and only if $\sigma=1/2$. Thus equal information weight and exact phase cancellation select the same point.

Third, the completed zeta function obeys a reflection symmetry about the line $\Ree s=1/2$. The anti-linear involution
\begin{equation}
J(s)=1-\overline{s}
\end{equation}
has fixed points exactly when $\Ree s=1/2$.

The mathematical programme is to determine whether these three facts are merely compatible or whether a canonical representation makes them logically dependent.

\section{Why the Berry phase is relevant}
The state $\ket{\psi(\sigma,\phi)}$ traces a latitude on the Bloch sphere as $\phi$ advances from $0$ to $2\pi$. At the balanced point $\sigma=1/2$, this latitude is the equator. A direct calculation gives the Berry phase
\begin{equation}
\gamma_B(\sigma)=-2\pi(1-\sigma)
\end{equation}
for the chosen gauge, hence
\begin{equation}
\gamma_B(1/2)=-\pi,
\qquad
\exp(i\gamma_B)=-1.
\end{equation}
This is the spinorial sign reversal associated with a full $2\pi$ cycle. The relevant content is not merely Euler's identity $e^{i\pi}=-1$; it is that the same phase arises as a geometric holonomy at the information-balanced equator. Berry's original formulation and Simon's line-bundle interpretation establish the standard geometric setting for this calculation \cite{Berry1984,Simon1983}.

\section{What would count as a proof}
A proof of the Riemann hypothesis within this programme would require at least one of the following:
\begin{enumerate}
\item a canonical map from each nontrivial zero to a normalized two-branch state for which zerohood is equivalent to exact destructive interference;
\item a self-adjoint operator whose spectral determinant is proportional to $\xi(1/2+iE)$;
\item a positive or totally positive kernel representation whose real-zero theorem applies to the Riemann $\Xi$ function;
\item a proof that the zeta functional equation acts locally on each zero state, rather than only globally by pairing distinct zeros.
\end{enumerate}
The present monograph completes neither item. It isolates the required implication in a form that can be attacked mathematically and computationally.

\section{Relation to established programmes}
The operator route is related to the Hilbert--Pólya idea and to spectral models developed by Berry and Keating, Connes, Sierra, and others \cite{BerryKeating1999,Connes1999,Sierra2007}. Those programmes ask whether the ordinates of the zeros arise as spectral data of a Hermitian or self-adjoint system. The distinctive proposal here is that the real part $1/2$ may be interpreted as a balance coordinate in a two-state information geometry and simultaneously as the equatorial location of a $\pi$ holonomy.

This interpretation is intended to sharpen the target, not to replace the operator problem. The spectral mechanism must still be constructed.

\chapter{The Completed Zeta Function and Its Fixed Axis}
\label{ch:zeta}

\section{Definition and analytic continuation}
For $\Ree s>1$, the Riemann zeta function is defined by the absolutely convergent Dirichlet series
\begin{equation}
\zeta(s)=\sum_{n=1}^{\infty}n^{-s}
\end{equation}
and by the Euler product
\begin{equation}
\zeta(s)=\prod_{p\ \mathrm{prime}}\left(1-p^{-s}\right)^{-1}.
\end{equation}
It extends meromorphically to the complex plane with a single simple pole at $s=1$. The nontrivial zeros lie in the critical strip $0<\Ree s<1$ and are symmetric about both the real axis and the critical line \cite{Riemann1859,Titchmarsh1986,DLMF25}.

A convenient entire completion is
\begin{equation}
\xi(s)=\frac12 s(s-1)\pi^{-s/2}\Gamma\!\left(\frac{s}{2}\right)\zeta(s).
\label{eq:xi-def}
\end{equation}
The factors outside $\zeta$ remove the pole and package the trivial zeros into the gamma factor. The completed function obeys
\begin{equation}
\xi(s)=\xi(1-s).
\label{eq:xi-fe}
\end{equation}
It also satisfies Schwarz reflection,
\begin{equation}
\xi(\overline{s})=\overline{\xi(s)},
\label{eq:xi-conj}
\end{equation}
because its defining data are real on the real axis and the analytic continuation respects conjugation.

\section{Linear and anti-linear symmetries}
There are two related transformations:
\begin{equation}
R(s)=1-s,
\qquad
J(s)=1-\overline{s}.
\end{equation}
Both are involutions:
\begin{equation}
R(R(s))=s,
\qquad
J(J(s))=s.
\end{equation}
The fixed points of $R$ are only the single complex point $s=1/2$, while the fixed points of $J$ form the entire vertical critical line.

\begin{theorem}[Fixed set of the zeta anti-involution]
For $s=\sigma+it$,
\begin{equation}
J(s)=s
\quad\Longleftrightarrow\quad
\sigma=\frac12.
\end{equation}
\end{theorem}

\begin{proof}
Since
\begin{equation}
J(\sigma+it)=1-(\sigma-it)=1-\sigma+it,
\end{equation}
the equation $J(s)=s$ is equivalent to $1-\sigma=\sigma$, hence $\sigma=1/2$. The imaginary coordinate remains arbitrary.
\end{proof}

Thus the critical line is not an arbitrary vertical line. It is the fixed locus of the reflection that combines the functional-equation complement $s\mapsto1-s$ with complex conjugation.

\section{Orbit structure of zeros}
If $\rho$ is a nontrivial zero, then the classical symmetries imply that
\begin{equation}
\rho,\quad \overline{\rho},\quad 1-\rho,\quad 1-\overline{\rho}
\end{equation}
are also zeros, with multiplicities preserved. If $\Ree\rho\neq1/2$ and $\Imm\rho\neq0$, these are generically four distinct points. If $\Ree\rho=1/2$, then $1-\overline{\rho}=\rho$, and the quartet collapses to the conjugate pair $\rho,\overline{\rho}$.

\begin{warning}
The functional equation guarantees orbit symmetry. It does not, by itself, require every zero to be a fixed point of $J$. Any proposed proof must supply a new reason why four-point orbits are impossible.
\end{warning}

This distinction is the logical center of the monograph. The information-spinor construction naturally selects fixed points, but a separate theorem is needed to show that all zero states belong to that fixed class.

\section{Centered coordinate}
Introduce
\begin{equation}
z=s-\frac12.
\end{equation}
Then the functional equation becomes an evenness relation. Define
\begin{equation}
\Xi_c(z)=\xi\!\left(\frac12+z\right).
\end{equation}
Equation \eqref{eq:xi-fe} gives
\begin{equation}
\Xi_c(z)=\Xi_c(-z).
\end{equation}
The critical line corresponds to purely imaginary $z=it$. The traditional real Riemann $\Xi$ function is
\begin{equation}
\Xi(t)=\xi\!\left(\frac12+it\right),
\end{equation}
which is real and even for real $t$.

The centered coordinate already has the form of a balance displacement. The real deviation
\begin{equation}
\delta=\sigma-\frac12
\end{equation}
changes sign under complement. The information model will identify this with the signed displacement from equal binary weight.

\section{The phase factor in the uncompleted functional equation}
One may write
\begin{equation}
\zeta(s)=\chi(s)\zeta(1-s),
\end{equation}
where one standard choice is
\begin{equation}
\chi(s)=2^s\pi^{s-1}\sin\!\left(\frac{\pi s}{2}\right)\Gamma(1-s).
\end{equation}
On the critical line $s=1/2+it$, the relation $1-s=\overline{s}$ and Schwarz reflection give
\begin{equation}
\zeta(s)=\chi(s)\overline{\zeta(s)}.
\end{equation}
Moreover $|\chi(1/2+it)|=1$. Therefore $\chi$ is a pure phase on the fixed axis. This established fact anticipates the later geometric statement: the critical line is where the dual branches are related by phase without a modulus imbalance.

Off the line, the duality factor generally changes modulus as well as phase. That observation motivates, but does not prove, the interpretation of $\sigma=1/2$ as the unitary or balanced duality axis.

\chapter{Binary Shannon Balance and Information Geometry}
\label{ch:shannon}

\section{Entropy in natural units}
For a Bernoulli distribution with probabilities $\sigma$ and $1-\sigma$, the Shannon entropy in nats is
\begin{equation}
\Hbin(\sigma)
=-\sigma\ln\sigma-(1-\sigma)\ln(1-\sigma),
\qquad 0\leq\sigma\leq1.
\label{eq:binary-entropy}
\end{equation}
The endpoint convention $0\ln0=0$ is understood. Shannon's original axiomatic derivation identifies the logarithm as the additive measure compatible with composition of choices; the base fixes the unit \cite{Shannon1948}. In base $e$, a balanced binary choice has entropy $\ln2$.

\begin{theorem}[Unique binary entropy maximum]
The function $\Hbin$ is strictly concave on $(0,1)$ and has the unique maximum
\begin{equation}
\Hbin\!\left(\frac12\right)=\ln2.
\end{equation}
\end{theorem}

\begin{proof}
Differentiation gives
\begin{equation}
\Hbin'(\sigma)=\ln\!\left(\frac{1-\sigma}{\sigma}\right),
\end{equation}
so the only stationary point satisfies $1-\sigma=\sigma$, hence $\sigma=1/2$. A second differentiation yields
\begin{equation}
\Hbin''(\sigma)
=-\frac1\sigma-\frac1{1-\sigma}
=-\frac{1}{\sigma(1-\sigma)}<0.
\end{equation}
Therefore the stationary point is the unique global maximum. Substitution gives
\begin{equation}
\Hbin(1/2)
=-2\left(\frac12\ln\frac12\right)=\ln2.
\end{equation}
\end{proof}

The number $\ln2$ is therefore not being inserted by hand at the balanced point. It is the exact natural-unit information of one maximally uncertain binary distinction.

\section{Complement symmetry}
Define the complement map
\begin{equation}
C(\sigma)=1-\sigma.
\end{equation}
Then
\begin{equation}
\Hbin(C(\sigma))=\Hbin(\sigma),
\end{equation}
and the unique fixed point of $C$ is $1/2$. This is formally parallel to the zeta complement $s\mapsto1-s$. The parallel becomes exact only after the model identification
\begin{assumption}[Real-part/probability identification]
\label{ass:sigma}
Within the information-spinor representation, identify the Bernoulli coordinate with the real part of the zeta argument:
\begin{equation}
\sigma=\Ree s.
\end{equation}
\end{assumption}

This assumption is natural because both coordinates lie in $(0,1)$ in their respective critical domains and both are acted on by the same complement $\sigma\mapsto1-\sigma$. Nevertheless, standard analytic number theory does not declare $\Ree s$ to be a probability. The identification is a modelling step and remains subject to justification by an operator or integral representation.

\section{Fisher metric of the Bernoulli family}
The Bernoulli family has Fisher information
\begin{equation}
g(\sigma)=\frac{1}{\sigma(1-\sigma)}.
\end{equation}
Thus the infinitesimal Fisher--Rao line element is
\begin{equation}
\dd \ell^2=\frac{\dd\sigma^2}{\sigma(1-\sigma)}.
\label{eq:fisher-line}
\end{equation}
The metric diverges at the deterministic boundaries and is symmetric about $1/2$. Introduce an angular coordinate $\theta$ by
\begin{equation}
\sigma=\cos^2\!\left(\frac{\theta}{2}\right)
=\frac{1+\cos\theta}{2}.
\label{eq:sigma-theta}
\end{equation}
Then
\begin{equation}
\dd\sigma=-\frac12\sin\theta\,\dd\theta,
\qquad
\sigma(1-\sigma)=\frac14\sin^2\theta,
\end{equation}
so \eqref{eq:fisher-line} reduces to
\begin{equation}
\dd \ell^2=\dd\theta^2.
\end{equation}
The Bernoulli manifold is therefore an interval with a natural angular coordinate. The balanced point $\sigma=1/2$ corresponds to $\theta=\pi/2$.

This is already the polar angle of the Bloch-sphere representation used in the next chapter. The information metric and the quantum-state geometry are thus joined by the square-root embedding rather than by metaphor.

\section{Square-root amplitude embedding}
Define
\begin{equation}
\boldsymbol{a}(\sigma)=
\begin{pmatrix}
\sqrt{\sigma}\\
\sqrt{1-\sigma}
\end{pmatrix}.
\end{equation}
Then $\|\boldsymbol{a}\|^2=1$, so the probability simplex is embedded into the unit circle in amplitude coordinates. The complement exchanges the two components. The balanced point becomes
\begin{equation}
\boldsymbol{a}(1/2)=\frac{1}{\sqrt2}
\begin{pmatrix}1\\1\end{pmatrix}.
\end{equation}

The square-root map is important because coherent cancellation acts on amplitudes, not directly on probabilities. A phase can rotate one amplitude relative to the other. The exact zero condition will therefore combine Shannon's equality of probabilities with a phase opposition of their square roots.

\section{Entropy deficit as a balance cost}
Define the entropy deficit
\begin{equation}
D_S(\sigma)=\ln2-\Hbin(\sigma)\geq0.
\end{equation}
Near $\sigma=1/2$, write $\sigma=1/2+\delta$. A Taylor expansion gives
\begin{equation}
\Hbin\!\left(\frac12+\delta\right)
=\ln2-2\delta^2-\frac43\delta^4+O(\delta^6),
\end{equation}
so
\begin{equation}
D_S(\sigma)=2\delta^2+\frac43\delta^4+O(\delta^6).
\end{equation}
The leading penalty for leaving the balance axis is quadratic. Any operator model that realizes $D_S$ as a positive energy or action would have a unique minimum at the critical line. This observation will later motivate a candidate confinement term, while remaining distinct from a proof that zeta zeros minimize it.

\begin{figure}[ht]
\centering
\includegraphics[width=0.76\textwidth]{binary_entropy.pdf}
\caption{Binary Shannon entropy in nats. The unique balance point is $\sigma=1/2$, where the entropy equals $\ln2$.}
\label{fig:entropy}
\end{figure}

\chapter{Bloch Geometry, Berry Holonomy, and the Spinor Sign}
\label{ch:bloch}

\section{Two-level state}
The square-root embedding becomes a normalized complex two-level state after introducing a relative phase:
\begin{equation}
\ket{\psi(\sigma,\phi)}
=\sqrt{\sigma}\ket{0}
+e^{i\phi}\sqrt{1-\sigma}\ket{1}.
\label{eq:qubit-state}
\end{equation}
Using \eqref{eq:sigma-theta}, this can be written
\begin{equation}
\ket{\psi(\theta,\phi)}
=\cos\!\left(\frac\theta2\right)\ket{0}
+e^{i\phi}\sin\!\left(\frac\theta2\right)\ket{1}.
\end{equation}
Up to an irrelevant global phase, this is the standard pure-qubit parametrization. Its Bloch vector is
\begin{equation}
\boldsymbol{n}=(\sin\theta\cos\phi,\sin\theta\sin\phi,\cos\theta).
\end{equation}
Since $\cos\theta=2\sigma-1$, the complement $\sigma\mapsto1-\sigma$ reflects the Bloch vector through the equatorial plane.

\begin{proposition}[Balanced point equals the Bloch equator]
The information-balanced condition $\sigma=1/2$ is equivalent to
\begin{equation}
n_z=0,
\end{equation}
so the full phase orbit lies on the Bloch equator.
\end{proposition}

\begin{proof}
From $n_z=\cos\theta=2\sigma-1$, one has $n_z=0$ if and only if $\sigma=1/2$.
\end{proof}

\section{Berry connection}
Let $\phi$ vary adiabatically around a full cycle while $\sigma$ is held fixed. Differentiate \eqref{eq:qubit-state}:
\begin{equation}
\partial_\phi\ket{\psi}
=i e^{i\phi}\sqrt{1-\sigma}\ket{1}.
\end{equation}
The Berry connection along the cycle is
\begin{equation}
\mathcal A_\phi
=i\braket{\psi}{\partial_\phi\psi}
=i\left[i(1-\sigma)\right]
=-(1-\sigma).
\end{equation}
Therefore the geometric phase is
\begin{equation}
\gamma_B(\sigma)
=\int_0^{2\pi}\mathcal A_\phi\,\dd\phi
=-2\pi(1-\sigma).
\label{eq:berry-sigma}
\end{equation}
A different gauge may shift $\gamma_B$ by an integer multiple of $2\pi$, but the holonomy
\begin{equation}
\mathcal U_B(\sigma)=e^{i\gamma_B(\sigma)}
\end{equation}
is gauge invariant.

The same result follows from the solid-angle formula. The latitude at polar angle $\theta$ encloses solid angle
\begin{equation}
\Omega=2\pi(1-\cos\theta),
\end{equation}
and a spin-$1/2$ state acquires
\begin{equation}
\gamma_B=-\frac{\Omega}{2}
=-\pi(1-\cos\theta)
=-2\pi(1-\sigma).
\end{equation}

\section{Equatorial holonomy}
\begin{theorem}[Information-balanced spinor holonomy]
For the state family \eqref{eq:qubit-state}, a full azimuthal cycle at the balanced point has
\begin{equation}
\gamma_B(1/2)=-\pi\pmod{2\pi},
\qquad
\mathcal U_B(1/2)=-1.
\end{equation}
\end{theorem}

\begin{proof}
Substituting $\sigma=1/2$ into \eqref{eq:berry-sigma} gives $\gamma_B=-\pi$. Euler's formula then gives $e^{-i\pi}=-1$.
\end{proof}

Thus the same point that maximizes binary entropy also produces the nontrivial sign holonomy of a spinor-like two-state cycle.

\section{Two pi and four pi}
For a spin-$1/2$ representation, a physical rotation by angle $\alpha$ about a unit axis $\boldsymbol{n}$ is represented by
\begin{equation}
U(\alpha)=\exp\!\left(-\frac{i\alpha}{2}\boldsymbol{n}\cdot\boldsymbol{\sigma}\right),
\end{equation}
where $\boldsymbol{\sigma}$ are the Pauli matrices. Therefore
\begin{equation}
U(2\pi)=-I,
\qquad
U(4\pi)=I.
\end{equation}
The factor $1/2$ is the representation weight responsible for the double cover $SU(2)\to SO(3)$. In the present construction, an analogous half-weight appears as the balance coordinate of the two complementary amplitudes.

It is essential to distinguish two claims:
\begin{enumerate}
\item \emph{Established:} spin-$1/2$ representations have a $-1$ sign under a $2\pi$ rotation.
\item \emph{Proposed identification:} the real coordinate of a zeta zero is represented by the population weight of a two-level state whose equatorial holonomy implements the zeta phase balance.
\end{enumerate}
The second is not implied merely by the first.

\section{Euler phase versus geometric phase}
The equation $e^{i\pi}=-1$ is algebraically universal. It does not by itself select $1/2$. The selection occurs because the Berry phase depends on the population coordinate:
\begin{equation}
\gamma_B(\sigma)=-2\pi(1-\sigma).
\end{equation}
Requiring the holonomy to equal $-1$ gives
\begin{equation}
-2\pi(1-\sigma)=\pi\pmod{2\pi},
\end{equation}
which yields
\begin{equation}
\sigma=\frac12\pmod1.
\end{equation}
In the probability interval $0\leq\sigma\leq1$, the nontrivial interior solution is $1/2$. The phase requirement therefore selects the same point as Shannon balance.

\begin{figure}[ht]
\centering
\includegraphics[width=0.66\textwidth]{bloch_equator.pdf}
\caption{Schematic Bloch geometry. Complement exchange reflects north and south populations, while $\sigma=1/2$ is the fixed equator.}
\label{fig:bloch}
\end{figure}

\chapter{The Exact Complementary-Cancellation Theorem}
\label{ch:cancellation}

\section{Scalar readout of a two-branch state}
Consider the coherent scalar readout
\begin{equation}
\mathcal A(\sigma,\phi)
=\sqrt{\sigma}+e^{i\phi}\sqrt{1-\sigma},
\qquad 0\leq\sigma\leq1.
\label{eq:amp}
\end{equation}
This is the overlap of the state \eqref{eq:qubit-state} with the unnormalized symmetric readout bra $\bra{0}+\bra{1}$. The associated intensity is
\begin{align}
\mathcal I(\sigma,\phi)
&=|\mathcal A(\sigma,\phi)|^2\\
&=\sigma+(1-\sigma)+2\sqrt{\sigma(1-\sigma)}\cos\phi\\
&=1+2\sqrt{\sigma(1-\sigma)}\cos\phi.
\label{eq:intensity}
\end{align}

\begin{theorem}[Unique exact cancellation]
\label{thm:cancellation}
For $0\leq\sigma\leq1$ and real $\phi$,
\begin{equation}
\mathcal A(\sigma,\phi)=0
\end{equation}
if and only if
\begin{equation}
\sigma=\frac12,
\qquad
\phi=(2k+1)\pi,
\quad k\in\mathbb Z.
\end{equation}
\end{theorem}

\begin{proof}
Suppose $\mathcal A=0$. Then
\begin{equation}
\sqrt{\sigma}=-e^{i\phi}\sqrt{1-\sigma}.
\end{equation}
Taking absolute values gives $\sqrt{\sigma}=\sqrt{1-\sigma}$, hence $\sigma=1/2$. Substituting this back yields $1+e^{i\phi}=0$, so $e^{i\phi}=-1$ and $\phi=(2k+1)\pi$. Conversely, those conditions clearly give $\mathcal A=0$.
\end{proof}

This theorem is the cleanest rigorous core of the information-spinor hypothesis. Exact zero formation in a normalized complementary pair requires both equal weights and phase opposition. Neither condition alone is sufficient.

\section{Alternative proof from the intensity bound}
Since
\begin{equation}
0\leq\sqrt{\sigma(1-\sigma)}\leq\frac12,
\end{equation}
with equality only at $\sigma=1/2$, equation \eqref{eq:intensity} implies
\begin{equation}
\mathcal I(\sigma,\phi)
\geq 1-2\sqrt{\sigma(1-\sigma)}\geq0.
\end{equation}
Equality in the first inequality requires $\cos\phi=-1$; equality in the second requires $\sigma=1/2$. This again proves Theorem \ref{thm:cancellation}.

\section{Factorization of the residual at phase pi}
At $\phi=\pi$,
\begin{equation}
\mathcal I(\sigma,\pi)
=\left(\sqrt{\sigma}-\sqrt{1-\sigma}\right)^2.
\end{equation}
Rationalizing gives
\begin{equation}
\mathcal I(\sigma,\pi)
=\frac{(2\sigma-1)^2}{\left(\sqrt{\sigma}+\sqrt{1-\sigma}\right)^2}.
\label{eq:intensity-delta}
\end{equation}
Hence the cancellation residual is an even, nonnegative function of the centered coordinate $\delta=\sigma-1/2$ and vanishes quadratically at the balance axis.

Near $\delta=0$,
\begin{equation}
\mathcal I\!\left(\frac12+\delta,\pi\right)
=2\delta^2+2\delta^4+O(\delta^6).
\end{equation}
The leading coefficient matches the leading Shannon entropy deficit derived in Chapter \ref{ch:shannon}:
\begin{equation}
\ln2-\Hbin\!\left(\frac12+\delta\right)
=2\delta^2+O(\delta^4).
\end{equation}
Thus, infinitesimally, informational imbalance and destructive-interference failure have the same quadratic cost.

\begin{proposition}[Local equivalence of balance costs]
As $\delta\to0$,
\begin{equation}
\frac{\mathcal I(1/2+\delta,\pi)}{\ln2-\Hbin(1/2+\delta)}\longrightarrow1.
\end{equation}
\end{proposition}

This local equality is stronger than the observation that both quantities have minima at $1/2$. It shows that, to second order, the Shannon information deficit is numerically identical to the loss of perfect phase cancellation in the normalized two-branch state.

\section{Geometric interpretation}
At $\sigma=1/2$, the state is
\begin{equation}
\ket{\psi(1/2,\phi)}
=\frac{1}{\sqrt2}\left(\ket0+e^{i\phi}\ket1\right).
\end{equation}
The points $\phi=0$ and $\phi=\pi$ are antipodal equatorial states. The cancellation readout at $\phi=\pi$ is orthogonal to the symmetric readout direction. Thus the zero is represented by a projective orthogonality event at the same equator that carries the $-1$ Berry holonomy.

\section{Generalized readout weights}
For completeness, consider
\begin{equation}
\mathcal A_{a,b}(\sigma,\phi)
=a\sqrt{\sigma}+b e^{i\phi}\sqrt{1-\sigma},
\end{equation}
with positive real $a,b$. Exact cancellation requires
\begin{equation}
\frac{\sigma}{1-\sigma}=\frac{b^2}{a^2},
\end{equation}
so
\begin{equation}
\sigma=\frac{b^2}{a^2+b^2}.
\end{equation}
Therefore $1/2$ is selected only by a symmetric readout $a=b$. The symmetry of the zeta functional equation is what motivates equal branch normalization. A canonical construction must derive, not merely choose, that symmetry.

\section{What this theorem does and does not establish}
The theorem proves the implication
\begin{equation}
\text{normalized complementary cancellation}
\Longrightarrow
\sigma=\frac12.
\end{equation}
It does not prove
\begin{equation}
\zeta(s)=0
\Longrightarrow
\text{normalized complementary cancellation}.
\end{equation}
The second implication is the open representation problem. Every later claim is organized around this logical boundary.

\begin{figure}[ht]
\centering
\includegraphics[width=0.76\textwidth]{destructive_interference.pdf}
\caption{At relative phase $\pi$, unequal complementary amplitudes leave a positive residual. Exact cancellation occurs only at $\sigma=1/2$.}
\label{fig:cancellation}
\end{figure}

\chapter{The Dirichlet Eta Bridge: Alternation, Parity, and ln 2}
\label{ch:eta}

\section{Alternating Dirichlet series}
The Dirichlet eta function is
\begin{equation}
\eta(s)=\sum_{n=1}^{\infty}\frac{(-1)^{n-1}}{n^s}.
\label{eq:eta-series}
\end{equation}
The alternating factor may be written
\begin{equation}
(-1)^{n-1}=e^{i\pi(n-1)},
\end{equation}
so the series carries an exact binary phase alternation. It converges for $\Ree s>0$ by the Dirichlet test and provides a useful continuation of the zeta relation into the critical strip.

Separate the zeta series into odd and even terms for $\Ree s>1$:
\begin{align}
\eta(s)
&=\sum_{n=1}^{\infty}n^{-s}-2\sum_{m=1}^{\infty}(2m)^{-s}\\
&=\zeta(s)-2^{1-s}\zeta(s)\\
&=\left(1-2^{1-s}\right)\zeta(s).
\label{eq:eta-zeta}
\end{align}
By analytic continuation the identity holds throughout the common domain.

\begin{theorem}[Eta-zeta factorization]
\begin{equation}
\eta(s)=\left(1-2^{1-s}\right)\zeta(s).
\end{equation}
Moreover,
\begin{equation}
\eta(1)=\ln2.
\end{equation}
\end{theorem}

\begin{proof}
The factorization was derived above where both series converge absolutely and extends analytically. At $s=1$, equation \eqref{eq:eta-series} becomes the alternating harmonic series, whose sum is $\ln2$.
\end{proof}

This is an exact point of contact among zeta, a binary sign phase, and the natural-unit binary information constant.

\section{Mellin transform of the fermionic kernel}
For $\Ree s>0$,
\begin{equation}
\int_0^\infty \frac{x^{s-1}}{e^x+1}\,\dd x
=\Gamma(s)\eta(s).
\label{eq:eta-mellin}
\end{equation}
Indeed,
\begin{equation}
\frac{1}{e^x+1}
=\sum_{n=1}^{\infty}(-1)^{n-1}e^{-nx},
\end{equation}
and termwise Mellin transformation yields \eqref{eq:eta-mellin}. The denominator $e^x+1$ is conventionally called fermionic because it appears in Fermi--Dirac statistics. This terminology does not by itself establish a physical fermion behind zeta zeros, but it reinforces the exact role of sign alternation.

At $s=1$,
\begin{equation}
\eta(1)
=\int_0^\infty\frac{\dd x}{e^x+1}
=\ln2.
\end{equation}
Thus $\ln2$ is simultaneously a Shannon binary entropy, an alternating harmonic sum, and the total weight of the simplest Fermi--Dirac kernel.

\section{Zeros introduced by the prefactor}
The factor
\begin{equation}
f(s)=1-2^{1-s}
\end{equation}
vanishes when
\begin{equation}
(1-s)\ln2=2\pi i k,
\qquad k\in\mathbb Z,
\end{equation}
that is,
\begin{equation}
s=1-\frac{2\pi i k}{\ln2}.
\end{equation}
These zeros lie on the line $\Ree s=1$, not in the open critical strip. Therefore within $0<\Ree s<1$, $\eta$ and $\zeta$ have the same zeros with the same multiplicities.

\begin{proposition}
For every $s$ with $0<\Ree s<1$,
\begin{equation}
\eta(s)=0\quad\Longleftrightarrow\quad\zeta(s)=0.
\end{equation}
\end{proposition}

The eta function is therefore a legitimate critical-strip representation of the nontrivial zero set.

\section{Why eta is not yet the missing proof}
The alternation $(-1)^{n-1}$ establishes a $0/\pi$ phase pattern over parity classes. However, the series contains infinitely many weighted terms, not an immediately normalized pair of complementary amplitudes. Grouping odd and even contributions gives
\begin{equation}
\eta(s)=O(s)-E(s),
\end{equation}
where in the absolute-convergence domain
\begin{equation}
O(s)=(1-2^{-s})\zeta(s),
\qquad
E(s)=2^{-s}\zeta(s).
\end{equation}
At a zeta zero, both analytically continued components vanish because both contain the common factor $\zeta(s)$. This factorization does not by itself show a nontrivial cancellation between two finite nonzero branch amplitudes.

The useful content of eta is therefore structural:
\begin{enumerate}
\item it places $\ln2$ and zeta in one exact identity;
\item it implements binary phase alternation;
\item it shares the nontrivial critical-strip zeros;
\item it suggests that a normalized parity decomposition may be possible.
\end{enumerate}
The missing step is a canonical renormalization or operator factorization that converts the infinite alternating sum into the two-branch cancellation form of Chapter \ref{ch:cancellation} without dividing by the very zero one seeks to explain.

\section{A proposed normalized parity problem}
Let $\mathcal H_o$ and $\mathcal H_e$ denote odd and even sectors of an appropriate Hilbert space. The research target is to construct states $\ket{O_s}$ and $\ket{E_s}$ and a bounded readout $L$ such that
\begin{equation}
\xi(s)=N(s)\,L\left(
\sqrt{\sigma}\ket{O_s}
+e^{i\Phi(s)}\sqrt{1-\sigma}\ket{E_s}
\right),
\end{equation}
where $N(s)$ is nonzero in the critical strip, complement exchanges the sectors, and $L$ has exact cancellation only at equal weights and phase $\pi$. This is not asserted to exist; it is a concrete formulation of the missing representation theorem.

\chapter{Critical-Line Phase Normalization and the Hardy Z Function}
\label{ch:hardy}

\section{Pure-phase duality on the fixed axis}
Write the functional equation as
\begin{equation}
\zeta(s)=\chi(s)\zeta(1-s).
\end{equation}
On $s=1/2+it$, one has $1-s=\overline{s}$, hence
\begin{equation}
\zeta\!\left(\frac12+it\right)
=\chi\!\left(\frac12+it\right)
\overline{\zeta\!\left(\frac12+it\right)}.
\label{eq:phase-duality}
\end{equation}
The modulus of $\chi$ is one on this line. Therefore there exists a real phase $\vartheta(t)$ such that
\begin{equation}
\chi\!\left(\frac12+it\right)=e^{-2i\vartheta(t)}.
\end{equation}
Define the Hardy function
\begin{equation}
Z(t)=e^{i\vartheta(t)}\zeta\!\left(\frac12+it\right).
\label{eq:hardyZ}
\end{equation}
Equation \eqref{eq:phase-duality} implies $Z(t)=\overline{Z(t)}$, so $Z(t)$ is real for real $t$. Its real zeros are precisely the zeros of zeta on the critical line \cite{Titchmarsh1986,DLMF25}.

This is an established phase-normalization result. It shows that the fixed axis converts the functional-equation duality into a real observable after removing a known phase.

\section{Balance of modulus and phase}
The critical line is distinguished by
\begin{equation}
|\chi(1/2+it)|=1.
\end{equation}
Thus duality changes phase but not magnitude. Away from the line, one generally has $|\chi(s)|\neq1$. In a two-branch interpretation, this means that the functional equation relates branches of unequal norm away from the fixed axis, while the axis supports unitary phase matching.

This statement must be handled carefully. A complex equation can vanish through cancellation even when individual terms have unequal raw norms if other coefficients or remainders are present. The modulus-one property is therefore a strong motivation for a balanced-state model, not an independent proof of RH.

\section{Approximate functional equation}
A standard approximate functional equation decomposes zeta into a direct Dirichlet sum and a dual sum multiplied by $\chi(s)$, schematically
\begin{equation}
\zeta(s)
=\sum_{n\leq x}n^{-s}
+\chi(s)\sum_{n\leq y}n^{s-1}
+R(s;x,y),
\qquad xy\asymp\frac{|t|}{2\pi}.
\label{eq:afe}
\end{equation}
On the critical line and with symmetric truncation $x\approx y$, the two sums are conjugate after phase normalization. This is the analytic-number-theory analogue of two complementary branches with matched weight.

The Riemann--Siegel formula sharpens this decomposition and is the basis of efficient evaluation of $Z(t)$. The observed zeros occur when the phase-normalized real quantity changes sign or reaches zero. The information-spinor programme proposes that the exact completion of this dual-sum structure may possess a projective two-channel interpretation.

\section{Riemann Xi as a real Fourier transform}
The completed function on the critical line admits a cosine-transform representation of the form
\begin{equation}
\Xi(t)=\int_0^\infty \Phi(u)\cos(tu)\,\dd u,
\label{eq:xi-fourier}
\end{equation}
for a rapidly decaying real kernel $\Phi$. Hence $\Xi$ is real and even. The Riemann hypothesis is equivalent to all zeros of this entire function occurring at real $t$.

This representation suggests a second route to the missing theorem. Instead of building a two-dimensional state for each zero, one may seek a positivity or total-positivity property of $\Phi$ strong enough to enforce real zeros of its Fourier transform. Classical work surrounding the de Bruijn--Newman deformation illustrates how heat-flow and zero geometry can encode the problem \cite{deBruijn1950,Newman1976,RodgersTao2020}.

\section{Phase quantization at a zero}
At a simple critical-line zero $t_n$, the real function $Z(t)$ crosses zero. Between zeros it has a sign, and its phase-normalized zeta value is either aligned or anti-aligned with the real readout. The cancellation theorem suggests representing the zero crossing by a relative phase of $\pi$ between dual contributions.

A rigorous implementation would require an exact identity
\begin{equation}
Z(t)=\mathcal N(t)\left[
\sqrt{\sigma(t)}F_+(t)
+e^{i\phi(t)}\sqrt{1-\sigma(t)}F_-(t)
\right]
\end{equation}
with $F_-$ canonically conjugate or complementary to $F_+$ and with nonvanishing normalization. On the critical line $\sigma=1/2$ would be fixed, and zerohood would reduce to phase opposition. The approximate functional equation provides a template but not yet such an exact normalized identity.

\begin{figure}[ht]
\centering
\includegraphics[width=0.84\textwidth]{hardy_z_zeros.pdf}
\caption{Numerical Hardy $Z(t)$ and the first ten tabulated critical-line zeros. This reproduces known data; it is not evidence that no off-line zeros exist.}
\label{fig:hardy}
\end{figure}

\chapter{The Conditional Information-Spinor Critical-Axis Theorem}
\label{ch:conditional}

\section{Abstract cancellation representation}
We now state the exact conditional result. Let $\rho=\sigma+it$ be a nontrivial zero of $\xi$. Suppose there exists a canonical normalized representation of its zero condition by two complementary branches. The word \emph{canonical} means that the representation is derived from zeta data and does not contain parameters fitted separately to each zero.

\begin{assumption}[Complementary zero-state representation]
\label{ass:zero-state}
For each nontrivial zero $\rho$ there exist normalized branch vectors $\ket{0_\rho}$ and $\ket{1_\rho}$, a nonzero scalar normalization $N(\rho)$, and a readout functional $L_\rho$ symmetric under branch exchange such that
\begin{equation}
\xi(\rho)=N(\rho)L_\rho\left[
\sqrt{\sigma}\ket{0_\rho}
+e^{i\phi(\rho)}\sqrt{1-\sigma}\ket{1_\rho}
\right],
\label{eq:zero-rep}
\end{equation}
with the following properties:
\begin{enumerate}[label=(A\arabic*)]
\item $\sigma=\Ree\rho$;
\item $J:\rho\mapsto1-\overline\rho$ exchanges the two branches;
\item symmetry of $L_\rho$ reduces zerohood to cancellation of equal readout amplitudes;
\item the relative zero phase is $\phi(\rho)=\pi\pmod{2\pi}$;
\item $N(\rho)\neq0$ and neither branch readout vanishes separately at $\rho$.
\end{enumerate}
\end{assumption}

Condition (A5) blocks a tautological factorization in which each branch already contains the factor $\xi(\rho)=0$.

\begin{theorem}[Conditional critical-axis selection]
\label{thm:conditional}
Under Assumption \ref{ass:zero-state}, every nontrivial zero $\rho$ satisfies
\begin{equation}
\Ree\rho=\frac12.
\end{equation}
\end{theorem}

\begin{proof}
By (A3)--(A5), the zero condition is a genuine cancellation between two nonzero complementary amplitudes. By (A4), their relative phase is $\pi$. Symmetric readout then reduces the condition to the amplitude of Theorem \ref{thm:cancellation}:
\begin{equation}
\sqrt{\sigma}-\sqrt{1-\sigma}=0.
\end{equation}
Hence $\sigma=1/2$. By (A1), $\sigma=\Ree\rho$, giving the result.
\end{proof}

\begin{corollary}
If Assumption \ref{ass:zero-state} is established for every nontrivial zero, then the Riemann hypothesis follows.
\end{corollary}

This corollary is logically correct but does not lessen the difficulty: proving the assumption is the substantive open task.

\section{Equivalent fixed-point formulation}
The same result can be expressed without explicit amplitudes.

\begin{assumption}[Local self-duality]
Each nontrivial zero is individually invariant under the anti-linear zeta involution:
\begin{equation}
J(\rho)=\rho.
\end{equation}
\end{assumption}

Then the fixed-set theorem of Chapter \ref{ch:zeta} immediately gives $\Ree\rho=1/2$. However, this assumption is nearly equivalent to RH itself. The amplitude formulation is more useful only if the local self-duality can be derived from independent normalization, unitarity, or positivity principles.

\section{Information and holonomy version}
A more geometric set of sufficient conditions is:
\begin{enumerate}
\item zero states are normalized points of a two-state projective space;
\item complement acts as the exchange of populations;
\item analytic continuation around the relevant cycle has holonomy $-1$;
\item the zero readout is orthogonality to the symmetric state;
\item the representation is faithful and nondegenerate.
\end{enumerate}
The holonomy condition yields $\sigma=1/2$ modulo one, while the probability domain selects the interior value $1/2$. Shannon entropy at that point is then $\ln2$.

\begin{proposition}[Triple coincidence under the model]
Under Assumption \ref{ass:sigma}, the following conditions are equivalent within the two-state representation:
\begin{align}
\sigma&=\frac12,\\
\Hbin(\sigma)&=\ln2,\\
C(\sigma)&=\sigma,\\
n_z&=0,\\
\mathcal U_B(\sigma)&=-1,\\
\mathcal A(\sigma,\pi)&=0.
\end{align}
\end{proposition}

\begin{proof}
The equivalence of the first three statements follows from Chapter \ref{ch:shannon}; the Bloch condition from $n_z=2\sigma-1$; the holonomy from Chapter \ref{ch:bloch}; and the cancellation condition from Theorem \ref{thm:cancellation}.
\end{proof}

This proposition is a genuine unification inside the model. It says that maximum binary distinction, complement fixedness, equatorial geometry, spinor sign, and destructive interference are different coordinates of the same balanced state.

\section{Why off-axis quartets remain possible without the assumption}
Suppose $\rho=1/2+\delta+it$ with $\delta\neq0$ is a hypothetical zero. The zeta symmetries generate the partner $J(\rho)=1/2-\delta+it$. A global two-state construction could assign those two distinct points as complementary states rather than forcing either one to be self-dual. In that case the pair as a whole is balanced even though each member is not.

Therefore a proof must exclude \emph{pairwise balance without pointwise balance}. This is the exact obstruction. Possible exclusion mechanisms include:
\begin{enumerate}
\item a self-adjoint spectral operator whose eigenparameter is real;
\item a positive norm that would be violated by paired nonreal eigenvalues;
\item a local unitarity law requiring $|\chi(s)|=1$ at every zero;
\item a canonical readout for which branch norms cannot be externally rescaled;
\item a total-positivity theorem for the Xi kernel.
\end{enumerate}

\section{A no-circularity criterion}
Any future proof claiming to establish Assumption \ref{ass:zero-state} should pass the following test:
\begin{quote}
Could the proposed branch states and normalization be defined before knowing that $\Ree\rho=1/2$, and do they remain nonsingular for a hypothetical off-axis zero?
\end{quote}
If the construction explicitly inserts $\sigma=1/2$, uses a list of critical-line zeros as input, or divides by a factor vanishing at the target zero, it is circular.

\chapter{Operator, Kernel, and Spectral Realization Programme}
\label{ch:operator}

\section{Hilbert--Pólya criterion}
A standard sufficient route to the Riemann hypothesis is to construct a self-adjoint operator $\widehat H$ whose spectrum consists of the ordinates of the nontrivial zeros. In its cleanest form, one seeks a relation such as
\begin{equation}
D(E)=C(E)\,\xi\!\left(\frac12+iE\right),
\label{eq:spectral-det}
\end{equation}
where $D(E)$ is a spectral determinant of $\widehat H$ and $C(E)$ is nonzero in the relevant domain. If $\widehat H$ is self-adjoint, its eigenvalues $E$ are real. The corresponding zeros then have real part $1/2$.

The principal difficulty is not writing an operator with a spectrum fitted to a pre-existing list of zeros. It is constructing a natural operator whose domain, boundary conditions, self-adjointness, and determinant are derived independently from arithmetic or geometry. Berry and Keating related the smooth zero-counting function to the semiclassical Hamiltonian $H=xp$, while Connes developed an absorption-spectrum interpretation through noncommutative geometry \cite{BerryKeating1999,Connes1999}. These are structural precedents rather than completed proofs.

\section{Information-spinor operator ansatz}
The present framework suggests a tensor-product Hilbert space
\begin{equation}
\mathcal H=\mathcal H_{\mathrm{arith}}\otimes\mathbb C^2,
\end{equation}
where the two-dimensional factor carries the complement degree of freedom. Let $\tau_x,\tau_y,\tau_z$ denote Pauli matrices acting on the complement factor. A schematic operator may be written
\begin{equation}
\widehat{\mathcal D}
=\widehat H_{\mathrm{arith}}\otimes I
+\widehat V\otimes\tau_x
+\widehat W\otimes\tau_y
+\widehat M\otimes\tau_z.
\label{eq:dirac-ansatz}
\end{equation}
The exchange involution acts as $\tau_x$ combined with complex conjugation. The balance axis corresponds to vanishing expectation of $\tau_z$:
\begin{equation}
\langle\tau_z\rangle=2\sigma-1=0.
\end{equation}
A chiral or particle-hole type symmetry could force eigenstates at a zero to have equal complement populations.

This is only an architecture. To become useful, the arithmetic operators must be specified and the determinant or scattering matrix must be shown to reproduce $\xi$.

\section{Positive balance penalty}
The Shannon deficit defines a nonnegative scalar potential
\begin{equation}
V_S(\sigma)=\ln2-\Hbin(\sigma).
\end{equation}
Likewise, the phase-$\pi$ cancellation residual is
\begin{equation}
V_I(\sigma)=\left(\sqrt\sigma-\sqrt{1-\sigma}\right)^2.
\end{equation}
Both vanish uniquely at $1/2$ and agree to quadratic order. An effective transverse Hamiltonian could contain
\begin{equation}
\widehat H_\perp
=-\alpha\frac{\dd^2}{\dd\sigma^2}
+\beta V_S(\sigma)
+\gamma V_I(\sigma),
\end{equation}
with suitable boundary conditions on $(0,1)$. Its ground state would be centered at $1/2$, but this alone would not prove that all zeta spectral states have zero transverse displacement. Excited transverse states generally exist. A supersymmetry, projection, topological index, or constraint would be needed to restrict the physical sector to the balanced mode.

\section{Holonomy as a boundary condition}
The Berry connection of Chapter \ref{ch:bloch} suggests imposing a monodromy condition
\begin{equation}
\Psi(\phi+2\pi)=-\Psi(\phi).
\label{eq:antiperiodic}
\end{equation}
Antiperiodic boundary conditions yield half-integer Fourier modes
\begin{equation}
\Psi_m(\phi)=e^{i(m+1/2)\phi},
\qquad m\in\mathbb Z.
\end{equation}
This is a precise mechanism by which a half shift appears spectrally. The challenge is to derive the zeta real part from this half-shift rather than merely reproducing a familiar spin structure.

One possible target is a transfer operator or Dirac-type operator whose Mellin exponent is shifted by the spin structure:
\begin{equation}
s=\frac12+iE.
\end{equation}
If the $1/2$ is a fixed spin-connection contribution and $E$ is the real spectrum of a self-adjoint generator, RH follows. The unresolved tasks are to identify the arithmetic evolution and prove the determinant identity \eqref{eq:spectral-det}.

\section{Mellin generator and dilation}
The classical generator of dilations is $xp$. Quantum mechanically, a symmetric realization is
\begin{equation}
\widehat D=\frac12(\widehat x\widehat p+\widehat p\widehat x)
=-i\left(x\frac{\dd}{\dd x}+\frac12\right)
\end{equation}
on an appropriate domain. The appearance of $1/2$ here is forced by symmetric ordering and by unitarity of dilations in $L^2(\mathbb R_+,\dd x)$. Indeed, the unitary dilation group acts as
\begin{equation}
(U_a f)(x)=e^{a/2}f(e^a x),
\end{equation}
and its generator contains the half-density term.

This provides a second exact origin of the half shift, independent of Shannon entropy. Mellin eigenfunctions behave as
\begin{equation}
f_E(x)=x^{-1/2+iE},
\end{equation}
up to conventions. Their exponent lies on the critical line. Berry--Keating type models exploit this dilation structure \cite{BerryKeating1999,Sierra2007}.

\begin{proposition}[Half-density origin of the critical line]
The unitary representation of dilations on $L^2(\mathbb R_+,\dd x)$ has generalized eigenfunctions with Mellin exponent $-1/2+iE$. Under the conventional zeta Mellin variable, this corresponds to a fixed real part $1/2$.
\end{proposition}

The proposition explains why self-adjoint dilation generators naturally live on the critical line. It still does not establish that their spectral boundary condition is the zeta function.

\section{Kernel criterion}
Another route starts with the cosine transform \eqref{eq:xi-fourier}. A sufficient result would be a theorem placing the kernel $\Phi$ in a class whose Fourier transform has only real zeros. Candidate concepts include Pólya frequency functions, total positivity, variation-diminishing transforms, and de Branges spaces.

The information interpretation suggests examining whether $\Phi$ can be normalized as a probability or amplitude kernel whose complement operation has a unique fixed point. The relevant test is not whether $\Phi$ is positive; positivity alone is insufficient to force real zeros of a Fourier transform. One needs a much stronger structural property.

\section{De Branges strategy}
De Branges theory gives conditions under which an entire function belongs to a Hermite--Biehler class and has controlled zero geometry. A possible target is to construct an entire function $E(z)$ such that
\begin{equation}
\Xi(z)=\frac12\left(E(z)+E^{\#}(z)\right),
\qquad
E^{\#}(z)=\overline{E(\overline z)},
\end{equation}
with
\begin{equation}
|E(z)|>|E^{\#}(z)|
\qquad \text{for }\Imm z>0.
\end{equation}
Such an inequality would encode a positivity or causality condition. In the two-branch language, $E$ and $E^{\#}$ are complementary amplitudes, and the real axis is their fixed boundary. Establishing the Hermite--Biehler inequality for an appropriate zeta-derived $E$ would be a substantive proof route rather than a restatement of symmetry.

\section{Minimal proof obligations}
A publishable operator claim must specify:
\begin{enumerate}
\item the Hilbert space and inner product;
\item the dense operator domain;
\item the adjoint and deficiency indices;
\item the self-adjoint extension, if any;
\item the boundary conditions and why they are arithmetic rather than fitted;
\item the determinant or scattering function;
\item proof that any prefactor multiplying $\xi$ is nonzero in the critical strip;
\item proof that no tabulated zero data were inserted into the operator.
\end{enumerate}
Without these items, an operator analogy remains heuristic.

\chapter{Correlation-kernel density criterion}\label{ch:correlation-kernel}

\section{A second analytic route from the Xi kernel}

The canonical Fourier-kernel representation introduced in the operator programme admits a second route that is independent of the Hermite--Biehler candidate audited in XF-2.  For the even Riemann kernel $\Phi$ and
\[
\Xi(z)=\int_{\mathbb R}\Phi(t)e^{izt}\,\dd t,
\]
define the second correlation function
\begin{equation}\label{eq:nu2-correlation}
\nu_2(t)
=
\int_{-\infty}^{\infty}
(t-2s)^2\Phi(t-s)\Phi(s)\,\dd s.
\end{equation}
For $0<|y|<1/2$, set
\begin{equation}\label{eq:phi2y-correlation}
\Phi_{2,y}(t)=\cosh(ty)\,\nu_2(t).
\end{equation}
These are the $n=2$ correlation objects used by Dimitrov and Xu in their Fourier--Wronskian criteria \cite{DimitrovXu2016}.

\section{Published RH-equivalent density criterion}

Dimitrov--Xu Theorem~1.1 supplies a standard external equivalence.  In the present notation,
\begin{equation}\label{eq:dx-density-rh}
\boxed{
\mathrm{RH}
\iff
\forall\,y\in\left(-\frac12,\frac12\right)\setminus\{0\},\quad
\overline{\operatorname{span}\mathcal T(\Phi_{2,y})}^{\,L^1(\mathbb R)}
=L^1(\mathbb R),
}
\end{equation}
where $\mathcal T(f)$ denotes the family of all translates of $f$.

The same work supplies an equivalent convolution-annihilator formulation: for every admissible $y$, the condition
\begin{equation}\label{eq:dx-annihilator}
\Phi_{2,y}*g=0,\qquad g\in L^\infty(\mathbb R),
\end{equation}
forces $g=0$ exactly on the RH surface \cite{DimitrovXu2016}.

Within the typed implication graph, the theorem itself has status \status{STANDARD}.  The global density condition and the equivalent annihilator condition have status \status{OPEN}; they are the mathematical premises that must be independently established before the standard equivalence can close the RH node.

\section{Wronskian bridge}

For $n=2$, the Fourier--Wronskian identity gives
\begin{equation}\label{eq:dx-wronskian}
W_2(\Xi;x)
=
\Xi(x)\Xi''(x)-\Xi'(x)^2
=
-\mathcal F[\nu_2](x),
\end{equation}
with the Fourier convention of \cite{DimitrovXu2016}.  Since $\Xi$ is real on the real $z$ axis, its first Laguerre quantity obeys
\begin{equation}\label{eq:laguerre-wronskian}
\boxed{
L_1[\Xi](x)
=
\Xi'(x)^2-\Xi(x)\Xi''(x)
=
-W_2(\Xi;x)
=
\mathcal F[\nu_2](x).
}
\end{equation}
This gives an exact consistency bridge between the correlation-kernel and derivative representations.

\section{Executable diagnostic layer}

The implementation in \texttt{src/critical\_axis/correlation\_kernel.py} evaluates $\nu_2$, $\Phi_{2,y}$, the Wronskian, and the Laguerre quantity with explicit finite truncation controls.  The validation suite checks:
\begin{enumerate}[label=(\roman*)]
\item evenness and positive reference values of the implemented kernel;
\item numerical evenness of $\nu_2$;
\item the exact $y\leftrightarrow-y$ symmetry of $\Phi_{2,y}$;
\item real-axis closure of Eq.~\eqref{eq:laguerre-wronskian};
\item fail-closed enforcement of $0<|y|<1/2$;
\item sampled Laguerre positivity as a \status{NUMERICAL\_DIAGNOSTIC} quantity.
\end{enumerate}

Finite samples and finite translate families carry diagnostic authority only.  The repository records this layer as
\[
\status{NUMERICAL\_DIAGNOSTIC}.
\]
The theorem-level target retains its full quantifiers over every admissible $y$, all translates, and the complete space $L^1(\mathbb R)$.

\section{Relation to the branch-cancellation programme}

The two active routes now share the same canonical Xi kernel but expose different closure obligations:
\begin{align}
\Phi
&\longrightarrow A_\pm
\longrightarrow \text{exact zero cancellation}
\longrightarrow \text{population/coordinate bridge},\\
\Phi
&\longrightarrow \nu_2
\longrightarrow \Phi_{2,y}
\longrightarrow \text{$L^1$ translation density or annihilator closure}.
\end{align}
The first route isolates a local zero-state representation problem.  The second route converts RH into a global harmonic-analysis completeness problem.  Their coexistence supplies a useful cross-audit: a future mechanism that closes the critical axis should be checked against both the local branch structure and the global correlation-kernel criterion.

\section{Current frontier}

The immediate XF-3 proof obligation is therefore
\begin{equation}\label{eq:xf3-frontier}
\boxed{
\forall\,0<|y|<\frac12:\quad
\overline{\operatorname{span}\mathcal T(\Phi_{2,y})}^{\,L^1(\mathbb R)}
=L^1(\mathbb R),
}
\end{equation}
or equivalently the global bounded-annihilator condition of Eq.~\eqref{eq:dx-annihilator}.  The solver records both as \status{OPEN\_RH\_EQUIVALENT\_CRITERION}; the Riemann-hypothesis node therefore retains status \status{OPEN}.

\chapter{Wiener--Laguerre scalar criterion}\label{ch:wiener-laguerre-scalar}

\section{Reduction of the density frontier}

Chapter~\ref{ch:correlation-kernel} expresses the Dimitrov--Xu critical-axis criterion as density of the translates of
\[
\Phi_{2,y}(t)=\cosh(ty)\nu_2(t),
\qquad
0<|y|<\frac12.
\]
Wiener's $L^1$ translation theorem allows this infinite-dimensional statement to be reduced to a pointwise condition on a Fourier transform.  The Wronskian identities of Dimitrov and Xu then identify that transform directly with derivatives of the completed Xi function \cite{DimitrovXu2016}.

For real $x$ and admissible $y$, define
\begin{equation}\label{eq:q-xi-definition}
Q_\Xi(x,y)
:=
|\Xi'(x+iy)|^2
-\Ree\!\left(
\Xi(x+iy)\overline{\Xi''(x+iy)}
\right).
\end{equation}

\section{Wronskian identity}

Write
\[
\Xi(x+iy)=\phi(x,y)+i\psi(x,y).
\]
The Cauchy--Riemann equations give
\begin{align}
Q_\Xi(x,y)
={}&\phi_x(x,y)^2-\phi(x,y)\phi_{xx}(x,y)\\
&+\psi_x(x,y)^2-\psi(x,y)\psi_{xx}(x,y).
\end{align}
With $W_2(f)=ff''-(f')^2$, this becomes
\begin{equation}\label{eq:q-xi-wronskian}
Q_\Xi(x,y)
=-W_2(\phi(\cdot,y);x)-W_2(\psi(\cdot,y);x).
\end{equation}
Dimitrov--Xu Theorem~1.3 / Theorem~2.5 and their Section~3 correlation identity yield
\begin{equation}\label{eq:q-xi-fourier}
\boxed{
Q_\Xi(x,y)=\mathcal F[\Phi_{2,y}](x).
}
\end{equation}
The repository records Eq.~\eqref{eq:q-xi-fourier} as \status{STANDARD\_PASS}.

\section{Tauberian equivalence}

For each fixed admissible $y$, Wiener's $L^1$ theorem states
\begin{equation}\label{eq:wiener-density-nonzero}
\overline{\operatorname{span}\mathcal T(\Phi_{2,y})}^{L^1(\mathbb R)}
=L^1(\mathbb R)
\iff
\mathcal F[\Phi_{2,y}](x)\neq0
\quad\forall x\in\mathbb R.
\end{equation}
The Dimitrov--Xu proof also gives strict positivity at the origin,
\begin{equation}\label{eq:q-origin-positive}
Q_\Xi(0,y)>0
\qquad
\left(0<|y|<\frac12\right).
\end{equation}
Because $Q_\Xi(\cdot,y)$ is continuous and real-valued on $\mathbb R$, a globally nonvanishing $Q_\Xi$ has the positive sign selected by Eq.~\eqref{eq:q-origin-positive}.  Therefore
\begin{equation}\label{eq:density-q-positive}
\boxed{
\overline{\operatorname{span}\mathcal T(\Phi_{2,y})}^{L^1}=L^1
\iff
Q_\Xi(x,y)>0\quad\forall x\in\mathbb R.
}
\end{equation}
Combining Eq.~\eqref{eq:density-q-positive} with Dimitrov--Xu Theorem~1.1 gives
\begin{theorem}[Wiener--Laguerre scalar RH criterion]\label{thm:wiener-laguerre-rh}
Under the standard Xi-kernel hypotheses used by Dimitrov and Xu,
\begin{equation}\label{eq:q-rh-equivalence}
\boxed{
\mathrm{RH}
\iff
\forall\,0<|y|<\frac12\ \forall x\in\mathbb R:\quad
Q_\Xi(x,y)>0.
}
\end{equation}
\end{theorem}

The equivalence in Theorem~\ref{thm:wiener-laguerre-rh} has status \status{STANDARD}.  The globally quantified sign condition has status
\[
\status{OPEN\_RH\_EQUIVALENT\_CRITERION}.
\]

\section{Executable diagnostic}

The function
\begin{center}
\texttt{xi\_wiener\_laguerre\_scalar(x, y)}
\end{center}
implements Eq.~\eqref{eq:q-xi-definition} and enforces the open domain $0<|y|<1/2$.  The tests verify $y\leftrightarrow-y$ symmetry, positive origin reference values, and fail-closed domain handling.

Finite samples retain \status{NUMERICAL\_DIAGNOSTIC} authority.  An analytic closure must carry both universal quantifiers in Eq.~\eqref{eq:q-rh-equivalence}.

\section{Certified falsification surface}

The scalar form provides a direct falsification target.  A certified pair $(x_*,y_*)$ with
\[
x_*\in\mathbb R,
\qquad
0<|y_*|<\frac12,
\qquad
Q_\Xi(x_*,y_*)\le0
\]
combined with Eq.~\eqref{eq:q-origin-positive} and continuity would force a real zero of $Q_\Xi(\cdot,y_*)=\mathcal F[\Phi_{2,y_*}]$.  The corresponding Wiener density condition would fail at that $y_*$.

This makes the next research frontier sharply local in representation while globally quantified in domain:
\begin{equation}\label{eq:xf4-frontier}
\boxed{
Q_\Xi(x,y)>0
\quad
\text{for every }x\in\mathbb R,
\quad0<|y|<\frac12.
}
\end{equation}
The next analytic task is to seek a structural representation of $Q_\Xi$ that certifies this sign over the entire strip.

\chapter{Nonlocal transverse curvature}\label{ch:nonlocal-curvature}

\section{From the Wiener--Laguerre scalar to geometry}

Chapter~\ref{ch:wiener-laguerre-scalar} isolates the RH-equivalent scalar condition
\begin{equation}\label{eq:xf5-q-frontier}
Q_\Xi(x,y)>0,
\qquad
x\in\mathbb R,
\qquad
0<|y|<\frac12.
\end{equation}
The next step is an exact geometric reformulation of the same scalar.  Set
\[
F_x(y):=|\Xi(x+iy)|^2.
\]
Since $\Xi$ is real-entire,
\[
\overline{\Xi(x+iy)}=\Xi(x-iy),
\]
and direct differentiation yields
\begin{equation}\label{eq:xf5-curvature-identity}
\boxed{
\frac12 F_x''(y)
=
|\Xi'(x+iy)|^2
-\Ree\!\left(
\Xi(x+iy)\overline{\Xi''(x+iy)}
\right)
=Q_\Xi(x,y).
}
\end{equation}
Equation~\eqref{eq:xf5-curvature-identity} has status \status{EXACT}.  Therefore the XF-4 global sign frontier is equivalent to strict transverse convexity,
\begin{equation}\label{eq:xf5-convexity-frontier}
\boxed{
F_x''(y)>0
\quad
\text{for every }x\in\mathbb R,
\quad0<|y|<\frac12.
}
\end{equation}
The universal statement in Eq.~\eqref{eq:xf5-convexity-frontier} retains
\[
\status{OPEN\_RH\_EQUIVALENT\_CRITERION}.
\]

\section{Connection with the classical growth criterion}

The function $F_x$ is even in $y$, hence
\[
F_x'(0)=0.
\]
Strict convexity on $0<y<1/2$ therefore implies
\begin{equation}\label{eq:xf5-inner-growth}
F_x'(y)>0,
\qquad 0<y<\frac12.
\end{equation}
The outer region $y\ge 1/2$, corresponding to $\Ree s\ge1$, is controlled by the standard symmetric Hadamard/logarithmic-derivative argument.  Combining the inner and outer regions recovers the classical positivity/growth route associated with Hinkkanen and Lagarias \cite{Lagarias1999,Planat2026Theta}.

The solver records the same implication chain in four breakable stages:
\begin{enumerate}
\item \path{Q_Xi strict positivity};
\item equivalently, strict transverse convexity;
\item consequently, positive vertical growth;
\item with the outer-halfplane closure, the classical RH-equivalent growth criterion.
\end{enumerate}
The first globally quantified condition retains \status{OPEN}.  Thus XF-5 changes the representation of the proof obligation and preserves its evidential status.

\section{Differentiated theta-kernel representation}

The repository normalization is
\begin{equation}\label{eq:xf5-repo-fourier}
\Xi(z)=2\int_0^\infty\Phi(u)\cos(zu)\,\dd u.
\end{equation}
Define
\[
D(z)=\int_0^\infty\Phi(u)\cos(zu)\,\dd u=\frac12\Xi(z).
\]
In diagonal coordinates
\[
u=a+b,
\qquad
v=a-b,
\qquad
a>|b|,
\]
write
\[
M(a,b)=\Phi(a+b)\Phi(a-b).
\]
The symmetrized first-growth kernel can be represented as
\begin{equation}\label{eq:xf5-growth-kernel}
K_{x,y}(a,b)
=
\frac a2\cos(2xb)\sinh(2ya)
+
\frac b2\cos(2xa)\sinh(2yb).
\end{equation}
Differentiation gives the local curvature kernel
\begin{equation}\label{eq:xf5-curvature-kernel}
\boxed{
L_{x,y}(a,b)
:=\partial_y K_{x,y}(a,b)
=
a^2\cos(2xb)\cosh(2ya)
+
b^2\cos(2xa)\cosh(2yb).
}
\end{equation}
Under Eq.~\eqref{eq:xf5-repo-fourier}, the XF-4 scalar becomes
\begin{equation}\label{eq:xf5-nonlocal-q}
\boxed{
Q_\Xi(x,y)
=4\iint_{a>|b|}
M(a,b)L_{x,y}(a,b)\,\dd a\,\dd b.
}
\end{equation}
The differentiated-kernel identity and normalization bookkeeping have status \status{EXACT}.  The sign of the complete integral over the full admissible domain remains the active global obligation.

\section{Exact positive anchor at the central slice}

For $x=0$,
\begin{equation}\label{eq:xf5-xzero-kernel}
L_{0,y}(a,b)
=a^2\cosh(2ya)+b^2\cosh(2yb)>0
\end{equation}
throughout the interior domain $a>|b|$, $a>0$.  Together with the positive theta mass $M(a,b)$, this gives pointwise positivity of the XF-5 curvature integrand at the central slice.  It supplies a kernel-level version of the positive-sign anchor used in the Wiener argument of Chapter~\ref{ch:wiener-laguerre-scalar}.

\section{Nonlocality firewall}

Planat's 2026 theta-kernel analysis derives an exact longitudinal--transverse decomposition of the Xi growth derivative \cite{Planat2026Theta}.  The companion August 2026 preprint proves, for that phase-aligned decomposition, that the two sectors cancel to all algebraic orders and that universal independent block positivity fails: sufficiently far in the oscillatory regime at least one aligned block is negative \cite{Planat2026Nonlocal}.

The repository therefore records the following research firewall:
\begin{itemize}
\item Global correlated sign control --- \status{ACTIVE}.
\item Exact resummation preserving sector correlation --- \status{ACTIVE}.
\item Positive representation of the fully coupled integral --- \status{ACTIVE}.
\item Universal independent phase-aligned block positivity.\par\noindent Status: \status{LITERATURE\_NO\_GO}.
\item Finite block sampling --- \status{NUMERICAL\_DIAGNOSTIC}.
\end{itemize}
The no-go result is representation-specific.  It narrows the admissible strategy to globally correlated control of Eq.~\eqref{eq:xf5-nonlocal-q} or to a different exact representation that preserves the sign-carrying remainder.

\section{Executable audit surface}

The module \path{critical_axis.nonlocal_curvature} exposes
\begin{itemize}
\item \path{xi_transverse_curvature(x,y)}, implementing Eq.~\eqref{eq:xf5-curvature-identity};
\item \path{xi_transverse_curvature_direct(x,y)}, independently differentiating $|\Xi|^2$ for regression;
\item \path{theta_growth_kernel(x,y,a,b)};
\item \path{theta_curvature_kernel(x,y,a,b)}, implementing Eq.~\eqref{eq:xf5-curvature-kernel};
\item \path{theta_curvature_integrand}, evaluating one local contribution with fail-closed diagonal coordinates.
\end{itemize}

The corresponding tests lock the factor of two in Eq.~\eqref{eq:xf5-curvature-identity}, compare analytic and direct second derivatives, differentiate Eq.~\eqref{eq:xf5-growth-kernel} independently, verify the $x=0$ positive anchor, and reject invalid $a,b$ domains.

\section{XF-5 frontier}

The new proof surface is therefore the complete correlated curvature integral
\begin{equation}\label{eq:xf5-final-frontier}
\boxed{
\iint_{a>|b|}
M(a,b)L_{x,y}(a,b)\,\dd a\,\dd b>0
}
\end{equation}
for every real $x$ and every $0<|y|<1/2$.  Viable next steps are an analytic global lower bound, an exact correlation-preserving resummation, a positive-definite representation of the fully coupled integral, or a certified counterexample on the admissible domain.  The universal sign statement retains \status{OPEN\_RH\_EQUIVALENT\_CRITERION}.


\chapter{Symbolic and Numerical Audit}
\label{ch:audit}

\section{Purpose of computation}
Numerics cannot prove that every zeta zero lies on the critical line. They can, however, verify the implementation of exact identities, detect algebraic mistakes, and test whether proposed normalizations become singular or circular. The repository therefore includes deterministic tests and a machine-readable audit.

\section{Audited identities}
The test suite checks:
\begin{enumerate}
\item $\Hbin(1/2)=\ln2$;
\item $\Hbin'(1/2)=0$ and $\Hbin''(1/2)<0$;
\item $\mathcal A(1/2,\pi)=0$;
\item $\mathcal A(\sigma,\pi)\neq0$ for representative $\sigma\neq1/2$;
\item $e^{i\gamma_B(1/2)}=-1$;
\item $J(1/2+it)=1/2+it$;
\item $\xi(s)=\xi(1-s)$ at high precision for generic complex points;
\item $\eta(s)=(1-2^{1-s})\zeta(s)$;
\item $\eta(1)=\ln2$;
\item reproduction of the first tabulated nontrivial zeros through the numerical library.
\end{enumerate}

The final check is explicitly labeled as reproduction of known critical-line zeros. It does not search the entire critical strip and cannot exclude off-axis zeros.

\section{Numerical stability}
The completed function
\begin{equation}
\xi(s)=\frac12s(s-1)\pi^{-s/2}\Gamma(s/2)\zeta(s)
\end{equation}
contains factors of very different magnitudes at large $|t|$. Arbitrary precision arithmetic is used to reduce cancellation error. The audit stores residuals rather than converting all checks to binary pass/fail at machine precision.

For a claimed exact identity, the expected residual should decrease when the working precision increases. A residual that plateaus indicates either an implementation error or an ill-conditioned formula.

\section{Claim-aware output}
The audit report separates four statuses:
\begin{description}
\item[Technical identities: PASS.] The code reproduces the exact formulas.
\item[Numerical reproduction: PASS.] The first selected tabulated zeros and symmetries are recovered.
\item[Critical-axis explanation: CONDITIONAL THEOREM.] The cancellation mechanism is exact given the representation assumption.
\item[Riemann hypothesis: OPEN.] No exhaustive proof is claimed.
\end{description}

This separation is preserved in both JSON and Markdown reports so that later code changes cannot silently promote a conditional result.

\section{Prospective computational tests}
The following tests are proposed for the next research phase:
\begin{enumerate}
\item construct exact dual branches from the approximate functional equation and quantify their norm mismatch as a function of $\sigma$;
\item test candidate normalizations on synthetic off-axis points rather than only on known zeros;
\item verify that a proposed branch map remains finite at zeros;
\item search for a positive-definite Gram representation of finite Xi-kernel truncations;
\item compute deficiency indices for candidate dilation-spinor operators;
\item compare the spectral counting function with the Riemann--von Mangoldt formula without fitting individual ordinates;
\item freeze all free parameters before evaluating high zeros not used in development.
\end{enumerate}

\section{Reproducibility commands}
From the repository root:
\begin{verbatim}
python3 -m pip install -e .[test]
python3 scripts/make_figures.py
python3 scripts/run_audit.py
pytest -q
make monograph
\end{verbatim}
The resulting PDF, figures, and reports are deterministic up to timestamps and platform-level font metadata.

\chapter{Falsification Criteria and Research Roadmap}
\label{ch:roadmap}

\section{How the model can fail}
The information-spinor interpretation should be abandoned or substantially revised if any of the following occurs:
\begin{enumerate}
\item no noncircular normalized two-branch representation of $\xi$ can be constructed;
\item every candidate representation requires insertion of $\Ree s=1/2$ before the zero condition is evaluated;
\item the complement branches can be rescaled arbitrarily, making the equal-weight conclusion gauge dependent;
\item a proposed operator is not essentially self-adjoint and no natural self-adjoint extension exists;
\item the spectral determinant reproduces only the average zero density, not the fluctuation spectrum;
\item the determinant identity contains an unanalysed prefactor with zeros in the critical strip;
\item numerical agreement disappears for frozen, out-of-sample ordinates;
\item the holonomy is a removable gauge artifact rather than a gauge-invariant cycle phase.
\end{enumerate}

These are not peripheral cautions. They define the empirical and mathematical content of the programme.

\section{Stage I: exact representation search}
The first objective is an exact zeta or Xi decomposition into complementary branches. Candidate starting points are:
\begin{enumerate}
\item the theta-function integral underlying the functional equation;
\item the Riemann--Siegel integral and approximate functional equation;
\item odd/even sectors of the eta representation;
\item Mellin transforms on $L^2(\mathbb R_+)$ with unitary dilation;
\item de Branges or Hermite--Biehler decompositions.
\end{enumerate}
Each construction must be evaluated by the no-circularity criterion of Chapter \ref{ch:conditional}.

\section{Stage II: derive equal branch normalization}
Even if two branches are found, the coefficient equality must be derived. Promising mechanisms include:
\begin{enumerate}
\item unitarity of the functional-equation scattering matrix on the fixed axis;
\item an exchange symmetry with a conserved norm;
\item a CPT-like antiunitary involution;
\item half-density normalization of the Mellin transform;
\item a projective metric that fixes the square-root amplitudes uniquely.
\end{enumerate}
A manually chosen normalization would weaken the result to a coordinate convention.

\section{Stage III: establish pointwise rather than pairwise self-duality}
This is the decisive stage. Global symmetry allows off-axis pairs. The project must identify a local principle that prevents a zero from being paired with a distinct partner. Candidate principles are:
\begin{enumerate}
\item self-adjointness of the spectral generator;
\item positivity of a Weil-type quadratic form;
\item a topological index that pins zero modes to the symmetry boundary;
\item a protected chiral zero mode in a two-component operator;
\item total positivity of the Xi kernel.
\end{enumerate}

\section{Stage IV: independent review}
Before any claim stronger than a conditional theorem, the work should undergo:
\begin{enumerate}
\item line-by-line review by an analytic number theorist;
\item operator-domain review by a functional analyst;
\item independent reproduction of the code and build;
\item adversarial search for hidden use of RH-equivalent assumptions;
\item explicit comparison with known failed proof patterns.
\end{enumerate}

\section{Publication ladder}
The work naturally separates into publishable units:
\begin{description}
\item[Paper A.] Exact equivalence of Shannon balance, complementary cancellation, Bloch equator, and Berry $-1$ holonomy.
\item[Paper B.] Dirichlet eta as a parity-phase bridge and the normalized parity representation problem.
\item[Paper C.] A candidate dilation-spinor operator with full domain and self-adjointness audit.
\item[Paper D.] Kernel positivity or de Branges formulation, if a nontrivial inequality is established.
\end{description}
Paper A can be mathematically complete without claiming RH. Papers C or D would contain the genuinely difficult number-theoretic advance.

\section{Versioning and frozen claims}
Version v0.1 freezes the following claims:
\begin{enumerate}
\item the exact bridge theorem inside the two-state model;
\item the conditional critical-axis theorem;
\item the explicit statement that the zero-state representation remains unproved;
\item the numerical audit protocol;
\item the operator and kernel proof obligations.
\end{enumerate}
Future versions must preserve this ledger append-only. A later correction must add a new record rather than rewrite the historical claim state.

\chapter{Conclusion}
\label{ch:conclusion}

The number $1/2$ is selected independently by three exact structures:
\begin{enumerate}
\item it is the unique fixed point of binary complement and the unique maximizer of Shannon entropy, where the information is $\ln2$;
\item it is the unique equal-population point at which two complementary amplitudes can cancel under a phase difference of $\pi$;
\item it is the Bloch equator at which a full azimuthal cycle carries Berry holonomy $-1$.
\end{enumerate}
The completed zeta function adds a fourth exact structure: $\Ree s=1/2$ is the fixed line of the anti-linear involution $s\mapsto1-\overline s$. The Dirichlet eta identity places zeta, binary alternation, and $\ln2$ in one exact analytic equation. The Hardy $Z$ function shows that on the critical line the functional-equation factor is a pure phase and the phase-normalized zeta value is real.

Version v0.2 adds a further \emph{compatibility structure}, not a new independent derivation of $1/2$. After centring with
\begin{equation}
u=s-\frac12,
\end{equation}
the zeta involution becomes $\widetilde J(u)=-\overline u$. Reciprocal inversion $\mathcal I(u)=1/u$ commutes with this involution and preserves the centred critical axis as a set. The separate radial compactification $u/(1+|u|)$ then realizes a finite disk with the balanced centre at $u=0$ and the asymptotic region approaching the boundary. These facts sharpen the geometry, but because the centring step already uses $1/2$, reciprocal invariance must not be counted as an independent proof of the critical line.

A second v0.2 refinement concerns the former probability-coordinate postulate. Under the explicit assumption that the strip-to-binary-simplex coordinate is affine and maps the strip endpoints to the simplex endpoints, $p=\Ree s$ is unique. This promotes the coordinate-theoretic component to a conditional theorem. It does not prove the stronger assertion that the canonical zero-state representation of $\xi$ must use that coordinate as a physical branch population.

The Aharonov--Bohm and Berry sectors can also be stated more precisely than before: both are realizations of $U(1)$ connection holonomy, and a dimensionless loop phase $\pi$ gives holonomy $-1$. The Hubble scale $L_H=c/H_0$ supplies a dimensionless radial normalization $H_0L/c$. Treating the local AB/Berry holonomy sector and the global Hubble normalization as two coordinates of one TIR potential-tension architecture remains a model postulate, not an established identity of gauge theory or cosmology. The broader holonomic connector $W_{ij}$ is only referenced here and remains formally delegated to TIR.

These results justify a precise interpretation:
\begin{quote}
The critical line is the unique information-balanced, phase-only, self-dual axis available to a normalized complementary two-state representation of the zeta functional equation, and the centred reciprocal geometry is exactly compatible with that axis.
\end{quote}

The interpretation is mathematically coherent and yields a valid conditional theorem. It does not yet prove that every nontrivial zero is represented by a pointwise self-dual cancellation state. Off-axis quartets remain compatible with the classical functional equation because symmetry can act between distinct zeros. Nor has it been shown that reciprocal inversion maps zeta zeros to zeta zeros.

The remaining problem has therefore been narrowed further. One must construct a canonical operator, kernel, connection, or normalized branch decomposition derived from zeta data for which the affine strip coordinate, the $U(1)$ loop phase, and the reciprocal symmetry are not imposed externally but arise as properties of the same faithful zero-state object. Such a construction must then exclude pairwise balance without local balance. If it is self-adjoint or positive in the required sense, the axis $1/2$ would no longer be an empirical location inserted into the spectral sector; it would be forced simultaneously by information balance, spinor holonomy, and spectral reality.

That is the central result and the central open debt of the information-spinor critical-axis programme after the v0.2 inverse-holonomy extension.


\appendix
\chapter{Detailed Derivations}
\label{app:derivations}

\section{Taylor expansion of binary entropy}
Let $\sigma=1/2+\delta$. Then
\begin{align}
\Hbin\!\left(\frac12+\delta\right)
={}&-\left(\frac12+\delta\right)
\ln\left(\frac12+\delta\right)
-\left(\frac12-\delta\right)
\ln\left(\frac12-\delta\right).
\end{align}
Using
\begin{equation}
\ln\left(\frac12\pm\delta\right)
=-\ln2+\ln(1\pm2\delta)
\end{equation}
and the power series for $\ln(1+x)$ gives
\begin{equation}
\Hbin\!\left(\frac12+\delta\right)
=\ln2-2\delta^2-\frac43\delta^4-\frac{32}{15}\delta^6+O(\delta^8).
\end{equation}
All odd powers cancel by complement symmetry.

\section{Taylor expansion of the cancellation residual}
At phase $\pi$,
\begin{equation}
\mathcal I
=\left(\sqrt{\frac12+\delta}-\sqrt{\frac12-\delta}\right)^2.
\end{equation}
Since
\begin{equation}
\mathcal I=1-2\sqrt{\frac14-\delta^2}
=1-\sqrt{1-4\delta^2},
\end{equation}
the binomial series yields
\begin{equation}
\mathcal I
=2\delta^2+2\delta^4+4\delta^6+10\delta^8+O(\delta^{10}).
\end{equation}
Thus the information deficit and interference residual coincide only to leading nonzero order; their higher-order coefficients differ.

\section{Fisher coordinate and Bloch angle}
For $\sigma=\cos^2(\theta/2)$,
\begin{equation}
\frac{\dd\sigma^2}{\sigma(1-\sigma)}
=\frac{(\frac14\sin^2\theta\,\dd\theta^2)}{\frac14\sin^2\theta}
=\dd\theta^2.
\end{equation}
The Fisher distance from the balanced point is therefore
\begin{equation}
d_F(\sigma,1/2)
=\left|\theta-\frac\pi2\right|
=\left|2\arccos\sqrt\sigma-\frac\pi2\right|.
\end{equation}

\section{Berry connection in two gauges}
The north-patch gauge used in the main text gives
\begin{equation}
\ket{\psi_N}
=\cos(\theta/2)\ket0+e^{i\phi}\sin(\theta/2)\ket1,
\end{equation}
with
\begin{equation}
\mathcal A_N=-\sin^2(\theta/2)\,\dd\phi
=-\frac{1-\cos\theta}{2}\dd\phi.
\end{equation}
A south-patch gauge may be chosen as
\begin{equation}
\ket{\psi_S}=e^{-i\phi}\cos(\theta/2)\ket0+\sin(\theta/2)\ket1,
\end{equation}
with
\begin{equation}
\mathcal A_S=\cos^2(\theta/2)\,\dd\phi
=\frac{1+\cos\theta}{2}\dd\phi.
\end{equation}
Their difference is
\begin{equation}
\mathcal A_S-\mathcal A_N=\dd\phi,
\end{equation}
so the integrated phases differ by $2\pi$ and the holonomy is identical.

\section{Modulus of the functional-equation factor}
On the critical line, the functional equation and conjugation give
\begin{equation}
\zeta(s)=\chi(s)\overline{\zeta(s)}.
\end{equation}
At a point where $\zeta(s)\neq0$, taking absolute values yields
\begin{equation}
|\zeta(s)|=|\chi(s)|\,|\zeta(s)|,
\end{equation}
so $|\chi(s)|=1$. The identity extends continuously through zeros. This argument uses the fixed-line relation $1-s=\overline s$ and therefore does not hold off the line.

\section{Half-density term in the dilation generator}
Define the unitary dilation group on $L^2(\mathbb R_+,\dd x)$ by
\begin{equation}
(U_a f)(x)=e^{a/2}f(e^a x).
\end{equation}
Unitarity follows from
\begin{align}
\|U_a f\|^2
&=\int_0^\infty e^a|f(e^a x)|^2\,\dd x\\
&=\int_0^\infty |f(y)|^2\,\dd y.
\end{align}
Differentiating at $a=0$ gives
\begin{equation}
\left.\frac{\dd}{\dd a}U_a f\right|_{a=0}
=\frac12 f+x f'.
\end{equation}
Therefore the self-adjoint generator convention $U_a=e^{-ia\widehat D}$ gives
\begin{equation}
\widehat D=-i\left(x\frac{\dd}{\dd x}+\frac12\right).
\end{equation}
The half term is required by the Jacobian of the unitary action.

\chapter{Repository and Reproducibility Protocol}
\label{app:repo}

\section{Directory structure}
\begin{verbatim}
information_spinor_critical_axis/
  README.md
  LICENSES.md
  CITATION.cff
  Makefile
  pyproject.toml
  monograph/
    main.tex
    chapters/
    references.tex
  src/critical_axis/
    core.py
  scripts/
    make_figures.py
    run_audit.py
  tests/
    test_core.py
  figures/
  reports/
  logs/changes.jsonl
  .github/workflows/compile.yml
\end{verbatim}

\section{Build environment}
The reference build uses Python 3.11 or later, \texttt{mpmath}, \texttt{numpy}, \texttt{matplotlib}, \texttt{pytest}, XeLaTeX, and \texttt{latexmk}. No network access is required after dependencies are installed.

\section{Append-only change record}
Every conceptual or code change receives a new JSON Lines entry with:
\begin{enumerate}
\item UTC timestamp;
\item version and change identifier;
\item scope;
\item rationale;
\item affected files;
\item claim impact;
\item validation performed.
\end{enumerate}
Historical entries are not edited or deleted. Corrections are appended as new entries.

\section{Reproduction levels}
\begin{description}
\item[Level 1: algebra.] Run the test suite and reproduce exact residuals.
\item[Level 2: figures.] Regenerate all plots from source.
\item[Level 3: publication.] Compile the XeLaTeX monograph and inspect the PDF.
\item[Level 4: independent implementation.] Reimplement the identities without importing the package.
\item[Level 5: adversarial audit.] Attempt to construct synthetic off-axis zeros or branch states and test whether the assumptions silently exclude them.
\end{description}

\section{Provenance rule}
Numerical zero ordinates produced by \texttt{mpmath.zetazero} are treated as reproduced reference data, not as a proof source. Any future spectral model must generate its eigenvalues without reading a zero table.

\backmatter
\chapter{References}
\begin{thebibliography}{99}

\bibitem{Riemann1859}
B. Riemann,
\emph{Ueber die Anzahl der Primzahlen unter einer gegebenen Grösse},
Monatsberichte der Berliner Akademie, 1859, pp. 671--680.

\bibitem{Shannon1948}
C. E. Shannon,
``A Mathematical Theory of Communication,''
\emph{Bell System Technical Journal}, vol. 27, pp. 379--423 and 623--656, 1948.
DOI: \href{https://doi.org/10.1002/j.1538-7305.1948.tb01338.x}{10.1002/j.1538-7305.1948.tb01338.x} and
\href{https://doi.org/10.1002/j.1538-7305.1948.tb00917.x}{10.1002/j.1538-7305.1948.tb00917.x}.

\bibitem{Titchmarsh1986}
E. C. Titchmarsh, revised by D. R. Heath-Brown,
\emph{The Theory of the Riemann Zeta-Function}, 2nd ed., Oxford University Press, 1986.

\bibitem{Edwards1974}
H. M. Edwards,
\emph{Riemann's Zeta Function}, Academic Press, 1974.

\bibitem{Bombieri2006}
E. Bombieri,
``The Riemann Hypothesis,'' in J. Carlson, A. Jaffe, and A. Wiles, eds.,
\emph{The Millennium Prize Problems}, Clay Mathematics Institute and American Mathematical Society, 2006.

\bibitem{DLMF25}
NIST Digital Library of Mathematical Functions,
Chapter 25, ``Zeta and Related Functions,'' especially Sections 25.4 and 25.10,
\url{https://dlmf.nist.gov/25}, accessed 29 July 2026.

\bibitem{DimitrovXu2016}
D. K. Dimitrov and Y. Xu,
``Wronskians of Fourier and Laplace Transforms,''
\emph{arXiv preprint arXiv:1606.05011}, 2016.
\url{https://arxiv.org/abs/1606.05011}.

\bibitem{AmariNagaoka2000}
S.-I. Amari and H. Nagaoka,
\emph{Methods of Information Geometry},
American Mathematical Society and Oxford University Press, 2000.

\bibitem{Berry1984}
M. V. Berry,
``Quantal Phase Factors Accompanying Adiabatic Changes,''
\emph{Proceedings of the Royal Society of London A}, vol. 392, no. 1802, pp. 45--57, 1984.
DOI: \href{https://doi.org/10.1098/rspa.1984.0023}{10.1098/rspa.1984.0023}.

\bibitem{Simon1983}
B. Simon,
``Holonomy, the Quantum Adiabatic Theorem, and Berry's Phase,''
\emph{Physical Review Letters}, vol. 51, no. 24, pp. 2167--2170, 1983.
DOI: \href{https://doi.org/10.1103/PhysRevLett.51.2167}{10.1103/PhysRevLett.51.2167}.

\bibitem{AharonovBohm1959}
Y. Aharonov and D. Bohm,
``Significance of Electromagnetic Potentials in the Quantum Theory,''
\emph{Physical Review}, vol. 115, no. 3, pp. 485--491, 1959.
DOI: \href{https://doi.org/10.1103/PhysRev.115.485}{10.1103/PhysRev.115.485}.

\bibitem{Hubble1929}
E. Hubble,
``A Relation between Distance and Radial Velocity among Extra-Galactic Nebulae,''
\emph{Proceedings of the National Academy of Sciences}, vol. 15, no. 3, pp. 168--173, 1929.
DOI: \href{https://doi.org/10.1073/pnas.15.3.168}{10.1073/pnas.15.3.168}.

\bibitem{BerryKeating1999}
M. V. Berry and J. P. Keating,
``The Riemann Zeros and Eigenvalue Asymptotics,''
\emph{SIAM Review}, vol. 41, no. 2, pp. 236--266, 1999.
DOI: \href{https://doi.org/10.1137/S0036144598347497}{10.1137/S0036144598347497}.

\bibitem{Connes1999}
A. Connes,
``Trace Formula in Noncommutative Geometry and the Zeros of the Riemann Zeta Function,''
\emph{Selecta Mathematica}, vol. 5, no. 1, pp. 29--106, 1999.
DOI: \href{https://doi.org/10.1007/s000290050042}{10.1007/s000290050042}.

\bibitem{Sierra2007}
G. Sierra,
``$H=xp$ with Interaction and the Riemann Zeros,''
\emph{Nuclear Physics B}, vol. 776, pp. 327--364, 2007.
DOI: \href{https://doi.org/10.1016/j.nuclphysb.2007.03.049}{10.1016/j.nuclphysb.2007.03.049}.

\bibitem{deBruijn1950}
N. G. de Bruijn,
``The Roots of Trigonometric Integrals,''
\emph{Duke Mathematical Journal}, vol. 17, pp. 197--226, 1950.

\bibitem{Newman1976}
C. M. Newman,
``Fourier Transforms with Only Real Zeros,''
\emph{Proceedings of the American Mathematical Society}, vol. 61, no. 2, pp. 245--251, 1976.

\bibitem{RodgersTao2020}
B. Rodgers and T. Tao,
``The de Bruijn--Newman Constant Is Non-negative,''
\emph{Forum of Mathematics, Pi}, vol. 8, e6, 2020.
DOI: \href{https://doi.org/10.1017/fmp.2020.6}{10.1017/fmp.2020.6}.

\bibitem{deBranges1968}
L. de Branges,
\emph{Hilbert Spaces of Entire Functions},
Prentice-Hall, 1968.

\bibitem{KatzSarnak1999}
N. M. Katz and P. Sarnak,
\emph{Random Matrices, Frobenius Eigenvalues, and Monodromy},
American Mathematical Society, 1999.

\bibitem{Montgomery1973}
H. L. Montgomery,
``The Pair Correlation of Zeros of the Zeta Function,''
in \emph{Analytic Number Theory}, Proceedings of Symposia in Pure Mathematics, vol. 24, pp. 181--193, 1973.

\end{thebibliography}

\printindex

\end{document}
